\documentclass[12pt, a4paper]{article}

\usepackage{graphicx}
\usepackage{setspace}
\usepackage{tikz}
\usetikzlibrary{positioning}

% STYLE
\usepackage[pdfborder={0 0 0}]{hyperref}
\usepackage{geometry}
\usepackage{numprint}

% BIBLIOGRAPHY
\usepackage{apacite} % after hyperref

\usepackage{epigraph}
\setlength\epigraphwidth{.8\textwidth}

% HYPOTHESES
% \usepackage{etoolbox,changepage}
\usepackage{ntheorem}
\newtheorem{req}{Research question}
\newtheorem{hyp}{Hypothesis}
% \AtBeginEnvironment{hyp}{\begin{adjustwidth}{3em}{3em}}
% \AtEndEnvironment{hyp}{\end{adjustwidth}}
\theoremseparator{:}

\title{The impact of web accessibility on the desire for avoidance}
\author{Tommaso Amici}

\begin{document}
\doublespacing
\pagenumbering{gobble}

\makeatletter
\begin{titlepage}
  \begin{center}
    \includegraphics[scale=0.5]{images/uva_logo.eps}

    \vspace{0.8cm}

    \huge{\textbf{\@title}}

    \vspace{0.8cm}

    \large{\@author}

    \normalsize{\today}

  \end{center}

  \vfill

  % \includegraphics[width=0.4\textwidth]{university}

  \normalsize

  Student number: 12306444

  MSc Business Administration - Digital Marketing

  Amsterdam Business School, University of Amsterdam

  Supervisor: Joris Demmers

  EBEC approval number: \todo{Guidelines 4.4.1 EBEC approval}

\end{titlepage}
\makeatother


\newpage

\pagenumbering{roman}

\begin{abstract}
  \#TODO
\end{abstract}
\newpage

\tableofcontents
\listoffigures
\listoftables
\newpage

\pagenumbering{arabic}

\section{Introduction}

\epigraph{The power of the Web is in its universality. Access by everyone regardless
  of disability is an essential aspect.}{Tim Berners-Lee, W3C Director and inventor
  of the World Wide Web}

\epigraph{Accessibility is the right thing to do.}{\citeA{valdes1999} for the United Nations}

\noindent
After more than twenty years, we have reached a tipping point: accessibility is not only
the right thing to do, but it is a legal requirement for public sector bodies in the EU
since the approval of the Web Accessibility Directive in 2016, and it may soon be a legal
requirement for private sector bodies should the European Accessibility Act be approved.

In \citeA{carter2001web}, a spokesperson for an electronics retailer when asked about
Internet users with disabilities is quoted as saying: ``that's not a market we've
thought about pursuing'' (p. 227). However, looking at some recent statistics paints a
striking picture: the Office for National Statistics estimates that, in 2020, 22\% of
internet users in the UK have self-assessed disabilities
\citeyear{ukinternetusers2020}, and according to different estimates 100 million people
in the EU and one billion people in the world have some form of disability
\cite{euWebAccessibility2021, whoDisability2021}.

Voluntarily excluding such a huge part of your customer base is shortsighted at best,
and illegal at worst. However, assuming they are managed rationally, businesses may be
operating under the impression that web accessibility does not affect purchase
intention, and as a consequence they will not dedicate resources to it.

In this paper, we will attempt to answer the following questions:

\begin{req}
  Does web accessibility have an effect on the desire for avoidance of disabled and non-disabled users?
\end{req}
\begin{req}
  Is accessibility a concern for online businesses, or are they justified in ignoring it, from an economic perspective?
\end{req}

The next section explains what web accessibility is, and why many websites fail to
properly support disabled users. Additionally, we describe the theoretical framework
which will be used to establish the relevancy of online complaints addressing web
accessibility.

\section{Theoretical background and hypotheses}

\subsection{Web accessibility}\label{sec:web-accessibility}

According to the Web Content Accessibility Guidelines (WCAG) 2.1, the objective of
developers and designers implementing accessible websites is to make content ``more
accessible to a wider range of people with disabilities, including accommodations for
blindness and low vision, deafness and hearing loss, limited movement, speech
disabilities, photosensitivity, and combinations of these, and some accommodation for
learning disabilities and cognitive limitations'' \cite[para. 1]{wcag21}.

It is important to note that disability can be temporary or situational: for example
people with a broken arm are temporarily disabled, while a parent holding a child is
situationally disabled, i.e., they cannot use two arms in that specific situation
\cite{henry2006}. This means that effectively any Internet user is potentially disabled
in a given moment. The respondents of a survey conducted in 2015 among web
accessibility specialists agreed that accessibility is applicable to everyone and that
it is highly related to usability \cite{yesilada2015}.

\citeA{carter2001web} nonetheless reports a second reason why accessibility is often
ignored by web developers and businesses: ``making the site accessible makes it
boring'' (p. 227). Designers are often convinced that building an accessible
site results in uninteresting and plain sites. Accessibility is thus considered useless
for non-disabled users, if not outright damaging to their experience, forcing them to
use a subpar product. On the contrary, \citeA{schmutz2016implementing} finds that
non-disabled users also benefit from a more accessible interface, and research by
\citeA{petrie2004} finds that accessibility does not constrain visual design.

The existing literature focuses primarily on the perceptions of businesses,
practitioners, and users, both disabled and non-disabled, highlighting how
accessibility is often ignored while it should not, both from a moral perspective and
from a practical perspective. However, a question that is not directly answered is
whether a lack of accessibility can lead to a loss of customers for an online business.

\citeA{morris2016} lightly touches upon this topic: analyzing the results of a survey
conducted on blind users of Twitter, they find that, as the platform becomes less
accessible with more and more users posting visual content rather than written content,
blind users' desire for avoidance increases because they simply cannot use the site in a
meaningful way anymore.

Because of all the aforementioned reasons, we predict that users complain online about
web accessibility, as it can be an afterthought, formally the hypothesis we intend to
test is:

\begin{hyp}
  Disabled users complain online about web accessibility
\end{hyp}
\begin{hyp}
  Non-disabled users complain online about web accessibility
\end{hyp}

However, despite evidence that companies have historically neglected web accessibility,
we predict that businesses nowadays promptly address complaints about web
accessibility, formally:

\begin{hyp}
  Complaints about web accessibility receive a response in a shorter time than other topics
\end{hyp}

In other words, we expect web accessibility to act as a moderator on the speed of the
responses from companies. In order to establish whether this is the case, we will draw
insights from the literature on electronic word of mouth (eWOM).

\subsection{Electronic word of mouth}\label{sec:ewom}

In this paper, we will adopt the definition of eWOM given by \citeA{hennigthurau2004}:
``any positive or negative statement made by potential, actual, or former customers
about a product or company, which is made available to a multitude of people and
institutions via the Internet'' (p. 38). In particular, we will restrict our focus to
social media, as they provide engagement metrics that the Internet in general does not.

When analyzing social media data, engagement metrics are useful for two reasons. First,
as a proxy of reach, as it is not possible to have exact statistics of how many people
saw a particular post or page, unless we are the publishers. Second, and this is
strictly related to the first reason, to control for the effect of activists, i.e.,
users who post often and who tend to be negative, as previous studies find empirically
that online reviews are dominated by them \cite{moe2012}. If we had to rely only on the
content that is published we may incorrectly attribute more importance to content that
was actually not widely shared or consumed.

Users that complain online tend to hold a grudge and although their desire for revenge
diminishes over time, their desire for avoidance does not and they will take their
business elsewhere \cite{tripp2011unhappy}. Moreover, if the complaint is public, other
users, so-called silent observers, may also be witnesses to the complaint and may start
avoiding a business as a consequence \cite{lopez2021}. We expect that a user's desire
for avoidance will be higher if the failure is related to accessibility rather than,
for example, shipping or price fairness, as the impossibility or reduced capability of
using a website is fundamental in choosing whether or not to keep using it. Thus, we
introduce another hypothesis:

\begin{hyp}
  Users' expressed desire for avoidance is relatively higher when complaining about web
  accessibility
\end{hyp}

When it comes to disabled users specifically, the literature finds they rely on online
reviews to bypass physical requirements, but while a recent empirical study in the
hotel industry finds that disabled guests' perception of risk is positively mitigated
by trustworthy reviews, it finds no significant effect on purchase decision
\cite{williams2006, zhang2021}. However, this should not be considered conclusive as
other studies conducted on the general population find that trustworthy online reviews
do have a positive effect on purchase intention \cite{chevalier2006, chakraborty2019,
  ventre2020}, and ultimately we are not concerned only with disabled users, because all
users benefit from accessible interfaces, as discussed in
\autoref{sec:web-accessibility}.

\section{Research method}

As mentioned in \autoref{sec:ewom}, we are restricting our data collection to social
media because they provide engagement metrics, more specifically we collected data from
Twitter. Among all the different social media sites, Twitter is more appropriate than
others because of its public nature, and because it allows users to both capture a
company's attention and to reach a wide audience with their complaining
\cite{lopez2021}.

Additionally, unlike research that focuses on online reviews, which will be usually
posted on a retailer's website (e.g., Amazon) or on websites specialized in reviews
(e.g., Yelp), this paper deals with meta complaints about the websites themselves and
about the customer experience, not about products customers have bought. For this
reason, we believe Twitter is more likely to be the place where to find meta
complaints.

\subsection{Data collection}

Through the Twitter API, \#XXX tweets posted between \#XXX date and \#XXX date were
collected. The content of each tweet has been categorized using a classifier trained on
a manually labeled subset of data.

\section{Results}

\section{Discussion}
\subsection{Theoretical contributions}
\subsection{Managerial implications}
\subsection{Limitations and future research}

\newpage
\bibliographystyle{apacite}
\bibliography{references}

\end{document}
